\section{Discussion}
\noindent \textbf{Limitations.}
Our study has several limitations. First, although \textsc{MalSkillsBench} is constructed from real-world incidents, it may not fully cover the diversity of malicious skills in the wild. Publicly reported and archived cases are inherently biased toward attacks that were discovered and preserved, and \textsc{MalSkillsBench} does not reflect the naturally imbalanced distributions of real-world registries. That said, we believe it still provides a realistic and reproducible foundation for comparison, especially given the current lack of public benchmarks.
Second, while \toolname supports broad multi-language parsing, deeper operand resolution currently depends on YASA~\cite{yasa-engine,wang2026yasa} and is limited to Python, JavaScript, Java, and Go. As a result, value-flow analysis may be less useful in other languages. Nevertheless, this limitation is not fundamental to \toolname, and future work can incorporate additional analyzers to further improve the prototype.
Third, as malicious-skill ecosystems evolve, attackers may adopt more implicit or adversarial strategies that reduce explicit security signals. Although \toolname is intentionally designed to avoid direct maliciousness judgment and instead reason over extracted sensitive operations, it may still be affected by AI-evasion strategies~\cite{checkpoint2025aievasion,melo2025aimalwarepromptinjection}. More broadly, malicious-skill detection remains an ongoing hide-and-seek between attackers and defenders, and we therefore view this work as a foundation for future research rather than a once-and-for-all solution.

\noindent \textbf{Implications.}
Our findings suggest three broader implications for securing the emerging agentic supply chain. First, defenses for traditional software supply chains cannot be directly transferred to skill ecosystems. Unlike traditional packages, the risk of a skill often emerges from interactions among heterogeneous artifacts, so the defenses should treat cross-artifact reasoning as a basic requirement.
Second, our results indicate that current protection mechanisms remain insufficient in practice. Although some registries have already integrated protections and the community has developed a range of scanning tools, our benchmark results and real-world findings show that these mechanisms still miss malicious skills or incur substantial false positives. 
Third, we believe that both the \textsc{MalSkillsBench} and the neuro-sysbolic design of \toolname can serve as a foundation for future research and inspire more effective defenses for these evolving threats.

